\subsection{Algoritmo Genético}
\label{subsec:genetico}

No Algoritmo Genético implementado, cada indivíduo representa uma rota completa de manutenção válida. Cada gene é a posição no agendamento, e o alelo é o identificador do ativo a ser visitado. Como a função $h$ calcula multas financeiras (logo, o objetivo é minimizar), utilizamos a seleção por torneio: sorteamos pequenos grupos da população e o indivíduo com o menor custo é escolhido para se reproduzir.

Usar a codificação por permutação impede o uso de cruzamentos simples (de 1 ou $k$ pontos), pois isso faria com que ativos fossem duplicados ou esquecidos. Para resolver isso, usamos o operador \textit{Order Crossover} (OX), que preserva a validade das rotas. Ele copia um trecho contínuo do primeiro pai e preenche o restante do cromossomo respeitando a ordem circular dos ativos do segundo pai. A Tabela \ref{tab:ag-cruzamento} mostra um exemplo passo a passo.

\begin{table}[H]
\centering
\caption{Exemplo passo a passo do cruzamento \textit{Order Crossover} (OX) adotado.}
\label{tab:ag-cruzamento}
\begin{tabular}{ll}
\toprule
Etapa                   & Permutação \\
\midrule
Pai 1                   & (3, \textbf{1, 5, 2}, 4) \\
Pai 2                   & (4, 2, 1, 5, 3) \\
Segmento herdado        & (--, \textbf{1, 5, 2}, --) \\
Preenchimento (Pai 2)   & (4, --, --, --, 3) \\
Filho                   & (4, 1, 5, 2, 3) \\
\bottomrule
\end{tabular}
\end{table}

\begin{table}[H]
\centering
\caption{Parâmetros do algoritmo genético por instância.}
\label{tab:ag-parametros}
\begin{tabular}{lccccc}
\toprule
Instância     & $N$ & $P_c$ & $P_m$ & Torneio & Elite \\
\midrule
pequena & 100 & 0.90 & 0.20 & 3 & 2 \\
intermedi\'aria & 100 & 0.90 & 0.20 & 3 & 2 \\
grande & 100 & 0.90 & 0.20 & 3 & 2 \\
\bottomrule
\end{tabular}
\end{table}

A mutação é feita aplicando o operador \textit{Swap} com uma probabilidade baixa ($P_m$). O elitismo também foi implementado para garantir que as 2 melhores rotas da população passem diretamente para a geração seguinte sem sofrer alterações. O algoritmo para imediatamente ao atingir o limite orçamental $B$, cortando gerações pela metade se necessário.

Monitorar a diversidade da população (número de cromossomos distintos, mostrado na Figura \ref{fig:ag-diversidade}) é crucial: se a população convergir rápido demais para soluções iguais, o algoritmo perde a capacidade de busca e vira apenas uma busca local cara.

\begin{figure}[H]
\centering
\figuraresultado{ag_diversidade.png}{Diversidade da população x geração (mediana das execuções), uma curva por instância.}
\caption{Evolução da diversidade da população do algoritmo genético.}
\label{fig:ag-diversidade}
\end{figure}